% generator_I08.tex -- Prop I.8 (Theor.): SSS congruence -- equal sides => equal angle.
% Geometry and text from jemmybutton (CC-BY-SA). Build: lualatex generator_I08.tex
\documentclass[12pt]{article}
\usepackage{geometry}\geometry{a5paper, margin=1.5cm}
\usepackage[lining]{ebgaramond}
\usepackage{amssymb}
\usepackage{byrne}
\def\mpPre{u := 1cm; textLabels := false;}

\begin{document}

\defineNewPicture[1/5]{%
  pair A, B, C, D, E, F, d;%
  A := (0, 0);%
  B := A shifted (-u, -4u);%
  C := A shifted (3/2u, -3u);%
  d := (0, -9/2u);%
  D := A shifted d;%
  E := B shifted d;%
  F := C shifted d;%
  byAngleDefine(F, D, E, byblack, SOLID_SECTOR);%
  byAngleDefine(C, A, B, byblack, SOLID_SECTOR);%
  draw byNamedAngleResized();%
  byLineDefine(A, B, byred,   SOLID_LINE, REGULAR_WIDTH);%
  byLineDefine(B, C, byblack, SOLID_LINE, REGULAR_WIDTH);%
  byLineDefine(C, A, byblue,  SOLID_LINE, REGULAR_WIDTH);%
  byLineDefine(D, E, byred,   SOLID_LINE, THIN_WIDTH);%
  byLineDefine(E, F, byblack, SOLID_LINE, THIN_WIDTH);%
  byLineDefine(F, D, byblue,  SOLID_LINE, THIN_WIDTH);%
  draw byNamedLineSeq(0)(CA,BC,AB);%
  draw byNamedLineSeq(0)(FD,EF,DE);%
  draw byLabelsOnPolygon(C, B, A)(ALL_LABELS, 0);%
  draw byLabelsOnPolygon(F, E, D)(ALL_LABELS, 0);%
}

% STATEMENT
\noindent
\begin{minipage}[t]{0.45\linewidth}
\centering\vspace{0pt}%
\drawCurrentPicture
\end{minipage}\hfill
\begin{minipage}[t]{0.52\linewidth}
\vspace{0pt}%
\textbf{Book~I.\enspace Prop.\ VIII. Theor.}

\medskip\noindent\itshape
If two triangles have two sides of the one respectively
equal to two sides of the other
($\drawUnitLine{CA} = \drawUnitLine{FD}$ and
$\drawUnitLine{AB} = \drawUnitLine{DE}$)
and also their bases ($\drawUnitLine{BC} = \drawUnitLine{EF}$),
equal; then the angles
(\drawFromCurrentPicture{%
  startAutoLabeling;%
  startGlobalRotation(-getAttribute("angle","Direction","A"));%
  draw byNamedAngleWithDummySides(A);%
  stopGlobalRotation;%
  stopAutoLabeling;%
} and
\drawFromCurrentPicture{%
  startAutoLabeling;%
  startGlobalRotation(-getAttribute("angle","Direction","D"));%
  draw byNamedAngleWithDummySides(D);%
  stopGlobalRotation;%
  stopAutoLabeling;%
}) contained by their equal sides are also equal.

\bigskip

% PROOF
\noindent\upshape
If the equal bases \drawUnitLine{BC} and \drawUnitLine{EF}
be conceived to be placed one upon the other, so that the
triangles shall lie at the same side of them, and that the
equal sides \drawUnitLine{AB} and \drawUnitLine{DE},
\drawUnitLine{CA} and \drawUnitLine{FD} be conterminous,
the vertex of the one must fall on the vertex of the other;
for to suppose them not coincident would contradict the last
proposition.

\smallskip
Therefore sides \drawUnitLine{AB} and \drawUnitLine{CA},
being coincident with \drawUnitLine{DE} and \drawUnitLine{FD},
$\therefore \drawAngle{A} = \drawAngle{D}$.

\begin{flushright}Q.\ E.\ D.\end{flushright}
\end{minipage}

\end{document}
